\chapter{ANALISIS DAN PERANCANGAN APLIKASI}

\section{Tinjauan Perusahaan}


Tuliskan di sini \textbf{Pedoman Penulisan:} Jika mahasiswa melakukan riset maka tinjauan perusahaan berisi sejarah mengenai institusi/perusahaan/organisasi tempat mahasiswa melakukan riset serta struktur yang ada di organisasi tersebut. Pada halaman lampiran dibuktikan dengan surat keterangan telah melakukan riset.


\subsection{Sejarah Perusahaan}
\subsection{Struktur Organisasi}

\section{Analisis Aplikasi}


Tuliskan di sini \textbf{Pedoman Penulisan:} Menjelaskan secara umum analisis yang dilakukan penulis sehingga menemukan permasalahan dan menjelaskan solusi atau permasalahan dengan mendesain software.


\subsection{Analisis Masalah}


Tuliskan di sini \textbf{Pedoman Penulisan:} Menjelaskan analisis yang dilakukan dalam penulisan sehingga menemukan dan mengidentifikasi permasalahan serta menjelaskan solusinya.


\subsection{Analisis Kebutuhan}


Tuliskan di sini \textbf{Pedoman Penulisan:} Setelah melakukan identifikasi dan mendefinisikan masalah, penulis merumuskan kebutuhan untuk menyelesaikan permasalahan.



\section{Rancangan Algoritma}


Tuliskan di sini \textbf{Pedoman Penulisan:} Pada sub bab ini dijelaskan atau diuraikan mengenai algoritma yang digunakan untuk menyelesaikan kasus (Contoh: pseudo code Algoritma Dijkstra).


